\documentclass[11pt]{article}
\usepackage[margin=1in]{geometry}
\usepackage{amsmath,amssymb,amsthm}
\usepackage{graphicx}
\usepackage[hidelinks]{hyperref}

\newtheorem{theorem}{Theorem}

\title{Literature-implied status: a subsymmetric direct-sum variant of Problem 8.1}
\author{Run \texttt{fa\_banach\_001}}
\date{\today}

\begin{document}
\maketitle

\section*{Status}

This packet records a \emph{literature-implied partial direct-sum variant}, not
a full solution of the source problem.  The source problem asks about arbitrary
bases with the factorization property in a complemented sum.  The later theorem
used here proves a positive result for $\ell^p$ direct sums under stronger,
subsymmetric structural hypotheses.

\section*{Source question}

The source paper is Lechner--Motakis--Mueller--Schlumprecht,
\emph{Strategically reproducible bases and the factorization property}
\cite{LMMS}.  In Section 8, Problem 8.1 asks:

\begin{center}
\includegraphics[width=\linewidth]{figures/open_problem_crop.png}
\end{center}

In words, if $X$ and $Y$ have bases with the factorization property, does the
union of those bases, in the right order, again have the factorization property
in the complemented sum of $X$ and $Y$?

\section*{Supporting theorem}

The supporting paper is Lechner, \emph{Subsymmetric bases have the
factorization property} \cite{LechnerSubsymmetric}.  Theorem 2.3 states:

\begin{center}
\includegraphics[width=\linewidth]{figures/supporting_theorem_crop.png}
\end{center}

The theorem applies to a sequence of Banach spaces $X_n$ paired with $Y_n$ and
coordinate vectors $(e_{n,j})_j,(f_{n,j})_j$ satisfying the paper's hypotheses
$(\mathrm{B1})$--$(\mathrm{B5})$ uniformly.  It further assumes that the
coordinate bases $(e_{n,j})_j$ are uniformly $(C_u,C_s)$-subsymmetric and
non-$\ell^1$-splicing.  Under these assumptions, the array
$(e_{n,j})_{n,j}$ has the factorization property in
$\ell^p((X_n))$ for every $1\leq p\leq\infty$.

\section*{Idea of the proof}

The source problem is hard because the factorization property need not behave
formally under arbitrary decompositions of a space.  Lechner's theorem supplies
a controlled setting where the decomposition is an $\ell^p$ direct sum and the
coordinate bases have enough symmetry to let a large diagonal operator be
blocked while preserving a large diagonal on an array subbasis.  The
factorization operators are then built on the array basis itself, giving
$I_{\ell^p((X_n))}=PTE$ for every operator $T$ with large diagonal.

\section*{Precise packet theorem}

\begin{theorem}
Let $1\leq p\leq\infty$.  In the setting of Lechner's Theorem 2.3, namely for
dual pairs $(X_n,Y_n,\langle\cdot,\cdot\rangle)$ satisfying
$(\mathrm{B1})$--$(\mathrm{B5})$ uniformly, and for uniformly subsymmetric,
non-$\ell^1$-splicing coordinate bases $(e_{n,j})_j$, the direct-sum array
basis $(e_{n,j})_{n,j}$ has the factorization property in $\ell^p((X_n))$.
Consequently, the $\ell^p$ direct-sum analogue of Problem 8.1 has a positive
answer for this subsymmetric class of bases.
\end{theorem}

\begin{proof}
This is exactly Theorem 2.3 of \cite{LechnerSubsymmetric}.  Given a bounded
operator
\[
  T:\ell^p((X_n))\longrightarrow \ell^p((X_n))
\]
with large diagonal relative to the array basis, i.e.
\[
  \inf_{n,j} |\langle T e_{n,j},f_{n,j}\rangle|>0,
\]
Theorem 2.3 gives bounded operators
\[
  E,P:\ell^p((X_n))\longrightarrow \ell^p((X_n))
\]
such that
\[
  I_{\ell^p((X_n))}=PTE.
\]
This is precisely the factorization property for the array basis in the
$\ell^p$ direct sum.

The relevance to Problem 8.1 is by direct identification: the array
$(e_{n,j})_{n,j}$ is the ordered union of the coordinate bases in the direct
sum model.  The identification is only partial because Problem 8.1 assumes
only that the individual bases have the factorization property and allows a
general complemented-sum situation, while Theorem 2.3 assumes the stronger
uniform subsymmetry and non-$\ell^1$-splicing hypotheses and works in
$\ell^p((X_n))$.
\end{proof}

\section*{Verification notes and limitations}

The packet uses a separate later paper and not a same-paper answer.  The later
paper does not explicitly state that it resolves Problem 8.1, so the result is
classified as \texttt{literature\_implied\_answer}, not
\texttt{literature\_already\_answered}.

The result should not be cited as a full answer to Problem 8.1.  It leaves open
the unrestricted complemented-sum question and the case of arbitrary bases with
the factorization property.  It also does not address the other open problems
in Section 8 of \cite{LMMS}.

\begin{thebibliography}{9}

\bibitem{LMMS}
R. Lechner, P. Motakis, P. F. X. Mueller, and T. Schlumprecht.
\newblock Strategically reproducible bases and the factorization property.
\newblock \emph{Israel Journal of Mathematics}, 238(1):13--60, 2020.
\newblock arXiv:1809.09423.

\bibitem{LechnerSubsymmetric}
R. Lechner.
\newblock Subsymmetric bases have the factorization property.
\newblock arXiv:2011.09915, 2020.

\end{thebibliography}

\end{document}
